\documentclass[carta,firma,sinnotas]{lafrox}
\usepackage{fvextra}
\usepackage{pdflscape}

\datoslafrox{
  documento = {Subdocumento 8: Anexos},
  subtitulo = {Anexos 8.A--8.G},
  alcance = {Gestión de riesgos de la propuesta LafroX},
  formulario = {T-16: Plan de riesgos},
  fecha = {},
  version = {0.1},
  nomenclatura = {LAFROX-Subdocumento8-Anexos},
}
\begin{document}
\portadalafrox
\preliminareslafrox
\setcounter{chapter}{7}

\chapter{Anexos del subdocumento 8}

En el presente documento se presentaran todos los anexos relacionados con la gestión de riesgos de la propuesta LafroX.

\section{Anexo 8.A : Análisis cualitativo de los riesgos}

La siguiente tabla muestra el listado completo de todos los riesgos, con su análisis relacionado con la probabilidad, impacto y exposición:

\begin{tablalafrox}{Valoración cualitativa vigente de riesgos. Fuente: T-16; bandas de prioridad SD8.}{tab:sd8-cuerpo-4}{>{\hsize=0.8\hsize\linewidth=\hsize\RaggedRight\arraybackslash}X >{\hsize=1.1\hsize\linewidth=\hsize\RaggedRight\arraybackslash}X >{\hsize=1.05\hsize\linewidth=\hsize\RaggedRight\arraybackslash}X >{\hsize=1.1\hsize\linewidth=\hsize\RaggedRight\arraybackslash}X >{\hsize=1.4\hsize\linewidth=\hsize\RaggedRight\arraybackslash}X}{\cab{ID} & \cab{Prob. ordinal} & \cab{Impacto ordinal} & \cab{Exposición ($P\times I$)} & \cab{Nivel de prioridad}}
R-001 & 3 & 2 & 6 & Media \\
R-002 & 3 & 1 & 3 & Baja \\
R-003 & 3 & 2 & 6 & Media \\
R-004 & 3 & 2 & 6 & Media \\
R-005 & 2 & 4 & 8 & Media \\
R-006 & 2 & 4 & 8 & Media \\
R-007 & 3 & 2 & 6 & Media \\
R-008 & 2 & 2 & 4 & Baja \\
R-009 & 2 & 2 & 4 & Baja \\
R-010 & 2 & 2 & 4 & Baja \\
R-011 & 2 & 4 & 8 & Media \\
R-012 & 3 & 4 & 12 & Alta \\
R-013 & 2 & 4 & 8 & Media \\
R-014 & 2 & 4 & 8 & Media \\
R-015 & 3 & 4 & 12 & Alta \\
R-016 & 2 & 4 & 8 & Media \\
R-017 & 3 & 2 & 6 & Media \\
R-018 & 3 & 2 & 6 & Media \\
R-019 & 3 & 1 & 3 & Baja \\
R-020 & 3 & 1 & 3 & Baja \\
R-021 & 3 & 2 & 6 & Media \\
R-022 & 3 & 2 & 6 & Media \\
R-023 & 2 & 2 & 4 & Baja \\
R-024 & 3 & 2 & 6 & Media \\
R-025 & 3 & 2 & 6 & Media \\
R-026 & 3 & 2 & 6 & Media \\
R-027 & 3 & 1 & 3 & Baja \\
R-028 & 2 & 2 & 4 & Baja \\
R-029 & 3 & 2 & 6 & Media \\
R-030 & 3 & 2 & 6 & Media \\
R-031 & 3 & 1 & 3 & Baja \\
R-032 & 2 & 2 & 4 & Baja \\
R-033 & 3 & 1 & 3 & Baja \\
R-034 & 2 & 2 & 4 & Baja \\
R-035 & 3 & 2 & 6 & Media \\
R-036 & 3 & 1 & 3 & Baja \\
R-037 & 3 & 1 & 3 & Baja \\
R-038 & 3 & 2 & 6 & Media \\
R-039 & 3 & 1 & 3 & Baja \\
R-040 & 3 & 1 & 3 & Baja \\
R-041 & 2 & 2 & 4 & Baja \\
R-042 & 3 & 2 & 6 & Media \\
R-043 & 2 & 2 & 4 & Baja \\
R-044 & 2 & 1 & 2 & Baja \\
R-045 & 3 & 2 & 6 & Media \\
\end{tablalafrox}


\clearpage
\section{Anexo 8.B: Insumos cuantitativos para un análisis Monte Carlo}

Este anexo organiza los datos cuantitativos disponibles en el CSV actualizado para preparar una simulación de exposición total en horas-hombre (HH). No presenta resultados de Monte Carlo. La ponderación de probabilidad de la Tabla 3 del CSV, que también se muestra junto al nivel ordinal en T-16, es distinta de la probabilidad inicial de ocurrencia de la Tabla 4. Para el análisis cuantitativo y el valor esperado se usa la probabilidad inicial de la Tabla 4, que es la empleada por la Tabla 6. Por ejemplo, para R-001, $45\%\times152{,}0\ \mathrm{HH}=68{,}400\ \mathrm{HH}$; la ponderación cualitativa de Tabla 3 para ese riesgo es $0{,}50$ y no se sustituye por la probabilidad inicial de $45\%$.

Las probabilidades iniciales son juicios de planificación sujetos a validación por los responsables; no representan frecuencias históricas demostradas. Los objetivos residuales se muestran solo como escenarios condicionados a evidencia y no como entradas base. Los impactos de HH son estimaciones puntuales; la suma de valores esperados de la fuente no equivale a un percentil de una simulación.

\begin{landscape}
\setlength{\textwidth}{\linewidth}
\begin{tablalafrox}{Insumos por riesgo para preparar la simulación. Fuente: CSV actualizado, Tablas 3--6 y 8.}{tab:monte-carlo-insumos}{G{0.06} G{0.11} G{0.12} G{0.14} G{0.14} G{0.15} Y}{\cab{ID} & \cab{Ponderación cualitativa\ (Tabla 3)} & \cab{P inicial QRA\ (Tabla 4)} & \cab{P residual objetivo$^{*}$} & \cab{Impacto condicional\ (HH)} & \cab{Valor esperado\ (HH)} & \cab{Mes estimado\ (Tabla 8)}}
R-001 & 50\% & 45\% & 20\% & 152,0 & 68,400 & 2 \\
R-002 & 50\% & 40\% & 20\% & 67,2 & 26,880 & 2 \\
R-003 & 50\% & 40\% & 20\% & 201,6 & 80,640 & 11 \\
R-004 & 50\% & 45\% & 25\% & 224,0 & 100,800 & 3 \\
R-005 & 30\% & 30\% & 10\% & 128,0 & 38,400 & 8 \\
R-006 & 30\% & 25\% & 10\% & 70,4 & 17,600 & 6 \\
R-007 & 50\% & 40\% & 20\% & 80,0 & 32,000 & 9 \\
R-008 & 30\% & 30\% & 10\% & 80,0 & 24,000 & 10 \\
R-009 & 30\% & 25\% & 10\% & 67,2 & 16,800 & 8 \\
R-010 & 30\% & 25\% & 10\% & 88,0 & 22,000 & 3 \\
R-011 & 30\% & 30\% & 10\% & 96,0 & 28,800 & 8 \\
R-012 & 50\% & 35\% & 10\% & 201,6 & 70,560 & 12 \\
R-013 & 30\% & 30\% & 10\% & 105,6 & 31,680 & 11 \\
R-014 & 30\% & 20\% & 10\% & 96,0 & 19,200 & 4 \\
R-015 & 50\% & 35\% & 10\% & 201,6 & 70,560 & 11 \\
R-016 & 30\% & 30\% & 10\% & 144,0 & 43,200 & 11 \\
R-017 & 50\% & 40\% & 20\% & 126,4 & 50,560 & 3 \\
R-018 & 50\% & 45\% & 25\% & 150,4 & 67,680 & 11 \\
R-019 & 50\% & 45\% & 25\% & 148,8 & 66,960 & 12 \\
R-020 & 50\% & 35\% & 20\% & 112,0 & 39,200 & 12 \\
R-021 & 50\% & 45\% & 20\% & 152,0 & 68,400 & 3 \\
R-022 & 50\% & 50\% & 25\% & 118,4 & 59,200 & 1 \\
R-023 & 30\% & 30\% & 10\% & 88,0 & 26,400 & 4 \\
R-024 & 50\% & 40\% & 20\% & 126,4 & 50,560 & 3 \\
R-025 & 50\% & 35\% & 20\% & 177,6 & 62,160 & 4 \\
R-026 & 50\% & 40\% & 25\% & 201,6 & 80,640 & 11 \\
R-027 & 50\% & 35\% & 20\% & 80,0 & 28,000 & 1 \\
R-028 & 30\% & 30\% & 15\% & 108,8 & 32,640 & 12 \\
R-029 & 50\% & 35\% & 20\% & 126,4 & 44,240 & 8 \\
R-030 & 50\% & 50\% & 35\% & 232,0 & 116,000 & 4 \\
R-031 & 50\% & 35\% & 20\% & 80,0 & 28,000 & 8 \\
R-032 & 30\% & 30\% & 15\% & 152,0 & 45,600 & 3 \\
R-033 & 50\% & 40\% & 20\% & 126,4 & 50,560 & 6 \\
R-034 & 30\% & 30\% & 15\% & 136,0 & 40,800 & 21 \\
R-035 & 50\% & 40\% & 25\% & 88,0 & 35,200 & 1 \\
R-036 & 50\% & 40\% & 20\% & 80,0 & 32,000 & 14 \\
R-037 & 50\% & 35\% & 20\% & 80,0 & 28,000 & 10 \\
R-038 & 50\% & 45\% & 20\% & 152,0 & 68,400 & 12 \\
R-039 & 50\% & 35\% & 20\% & 67,2 & 23,520 & 4 \\
R-040 & 50\% & 40\% & 20\% & 92,8 & 37,120 & 18 \\
R-041 & 30\% & 30\% & 15\% & 176,0 & 52,800 & 53 \\
R-042 & 50\% & 35\% & 20\% & 126,4 & 44,240 & 21 \\
R-043 & 30\% & 30\% & 10\% & 88,0 & 26,400 & 8 \\
R-044 & 30\% & 30\% & 10\% & 88,0 & 26,400 & 8 \\
R-045 & 50\% & 35\% & 15\% & 100,8 & 35,280 & 8 \\
\multicolumn{4}{r}{Totales de las Tablas 5--6} & 5.585,6 & 2.058,480 & -- \\
\end{tablalafrox}
\end{landscape}

\noindent\textit{$^{*}$ La probabilidad residual es un objetivo condicionado a las validaciones y respuestas del CSV; no se debe usar en el escenario base mientras esas condiciones no estén acreditadas. La columna de ponderación cualitativa se incluye para trazabilidad y no es la probabilidad de ocurrencia que usa la Tabla 6.}

\begin{tablalafrox}{Disponibilidad y límites de los datos para Monte Carlo. Fuente: revisión del CSV actualizado.}{tab:monte-carlo-preparacion}{F{0.20} Y F{0.18} Y}{\cab{Parámetro} & \cab{Información disponible} & \cab{Estado} & \cab{Qué falta para simular}}
Probabilidad inicial & Tabla 4: porcentaje inicial por riesgo y fundamento & Juicio de planificación, no frecuencia histórica & Validar con dueño; confirmar horizonte y recurrencia \\
Impacto en HH & Tablas 5--6: impacto puntual y valor esperado & Punto HH; sin distribución & Sustentar rango/distribución; separar componentes inciertos \\
Mes de ocurrencia & Tabla 8: mes estimado por riesgo & Punto, sin distribución temporal & Vincular a actividades/hitos y estimar dispersión \\
Probabilidad residual & Tabla 4: objetivo por riesgo y condición & Condicionada; eficacia pendiente & Acreditar la condición antes del escenario residual \\
Dependencias & El CSV no contiene matriz ni correlaciones & No disponible & Evaluar causas comunes; no inferir correlación por rama RBS \\
Cronograma & Tabla 15: referencias generales al plan/EDT y etapa & Sin red simulable & Incorporar tareas, duraciones/distribuciones, precedencias y mapeo riesgo--actividad \\
Costos & Consecuencias externas y costos de respuesta pendientes & Sin valorización integral & Completar CLP y costos de tratamiento para simular costo económico \\
\end{tablalafrox}

El total esperado de la Tabla 6 (2.058,480 HH) es analítico, no un resultado Monte Carlo. Simular HH con ocurrencia única, impactos puntuales y dependencias exige declarar esos supuestos; simular cronograma requiere duraciones inciertas, precedencias y mapeo riesgo--actividad. Al ejecutar el modelo, informar P50/P80/P90, reserva, excedencia, iteraciones y semilla.

\end{document}